\section{Lecture 5: Solid Abelian Groups}
\label{lec:5}

Scholze's fourth lecture isolates, inside $\CondAbl$, a good class of
``complete'' objects: the \emph{solid} abelian groups. The motivation is that
naive tensor products of condensed abelian groups are badly behaved --- their
underlying object is the algebraic (uncompleted) tensor product --- so one wants
a reflective subcategory of complete objects together with a completion functor,
so that completing a naive tensor product yields something reasonable. The
lecture gives a presentation of the theory of solid abelian groups that is
\emph{specific to the light setting} and quite different from the one in the
first \emph{Condensed Mathematics} notes \cite{clausen2019condensed}: solidity is
defined by the single condition that every null sequence be uniquely summable,
phrased through the internally projective object $P=\mathbb{Z}[\mathbb{N}\cup
\{\infty\}]/\mathbb{Z}[\infty]$ of \cref{lec:3}. From this one gets, essentially
for free, that $\Solid\subset\CondAbl$ is an abelian subcategory closed under all
the usual operations; that the reals are annihilated by solidification; and,
after a computation, that $\Solid$ has a single compact projective generator
$\prod_{\mathbb{N}}\mathbb{Z}$, recovering the structural facts stated in
\cref{lec:1,lec:4}. The lecture closes with the reading of solidity as the
ability to integrate against $\mathbb{Z}$-valued measures.

\subsection{Motivation: a class of complete objects}

\begin{remark}[Why one wants completeness]
\label{rmk:5:tensor-motivation}
Working inside $\CondAbl$, computations that start with reasonable examples and
only take limits and colimits stay reasonable. The trouble begins with
constructions like the tensor product. If one forms the tensor product of two
power-series algebras, the answer one wants is a \emph{completed} tensor product
--- morally a power-series algebra in both variables. But the naive tensor
product in condensed abelian groups has underlying object the ordinary algebraic
tensor product,
\begin{equation}
\label{eq:5:naive-tensor}
\bigl(M\otimes_{\CondAbl}N\bigr)(\ast)=M(\ast)\otimes_{\mathrm{Ab}}N(\ast),
\end{equation}
a nasty, essentially indescribable object. So tensor products of condensed
abelian groups are not nice, and the remedy is to single out a subclass of
\emph{complete} objects and to complete the naive tensor product into that
subclass, hoping for a reasonable answer.
\end{remark}

We want the notion of completeness to cut out an abelian category, closed enough
that cokernels of complete objects stay complete. Among the objects that should
count as complete are $\mathbb{Z}$, the $p$-adics $\mathbb{Z}_p$, and $\mathbb{Q}_p$;
and if the class is to be closed under extensions it must also contain non-Hausdorff
objects such as the quotient in
\begin{equation}
\label{eq:5:Zp-extension}
0\longrightarrow\mathbb{Z}\longrightarrow\mathbb{Z}_p\longrightarrow
\mathbb{Z}_p/\mathbb{Z}\longrightarrow0 .
\end{equation}
The quotient $\mathbb{Z}_p/\mathbb{Z}$ is a very non-Hausdorff thing, so one
\emph{cannot} phrase completeness by asking that convergent sequences have unique
limits. In the condensed setting there is no good abstract notion of a convergent
sequence other than one carrying a specified limit point, and even granting such
a notion one could not demand that limits exist, precisely because one wants to
allow these non-Hausdorff examples.

\begin{remark}[Nonarchimedean first, reals later]
\label{rmk:5:nonarchimedean-first}
Scholze recalled that when he and Clausen first thought about this, finding the
right notion of completeness was one of the key questions. They originally wanted
a single notion for which the real numbers, and ideally all locally compact
abelian groups, would be complete, while still allowing cokernels and staying
complete; they could not directly make that work. Forgetting the reals for the
moment and aiming only for a nonarchimedean theory, there is something that works
very cleanly. So it is difficult to find a notion under which $\mathbb{R}$ is
complete, but there is a theory that works well for all nonarchimedean examples;
the real numbers are recovered later, by a different story (the liquid theory).
\end{remark}

\begin{remark}[A presentation special to the light setting]
\label{rmk:5:light-presentation}
The presentation below is very different from the one in the first
\emph{Condensed Mathematics} lecture notes \cite{clausen2019condensed}, and it
really only works in the light setting. There, showing that solid abelian groups
form an abelian category was among the \emph{last} things one could prove; here it
will be essentially the \emph{first}.
\end{remark}

The guiding idea is a basic fact of nonarchimedean analysis: a sequence is
summable if and only if it is a null sequence (tends to $0$). We turn this into a
definition of completeness.

\begin{equation}
\label{eq:5:summable-idea}
\boxed{\;M\in\CondAbl \text{ is ``nonarchimedean complete'' if every null
sequence in } M \text{ is (uniquely) summable.}\;}
\end{equation}

\subsection{The definition via the free group on a null sequence}

To formalize \eqref{eq:5:summable-idea} we use the object that classifies null
sequences.

\begin{notation}[The free group on a null sequence]
\label{not:5:P}
Write
\begin{equation}
\label{eq:5:P}
P:=\mathbb{Z}[\mathbb{N}\cup\{\infty\}]\big/\mathbb{Z}[\infty]\ \in\ \CondAbl,
\end{equation}
the free condensed abelian group on a null sequence (the object $M$ of
\cref{def:3:null-sequence-group}). Maps out of $P$ are exactly null sequences in
the target. Recall from \cref{thm:3:three-properties}(3) that $P$ is
\emph{internally projective}.
\end{notation}

Internal projectivity says precisely that the internal Hom out of $P$ is exact.
Recall the internal Hom functor
\begin{equation}
\label{eq:5:internal-hom-functor}
\iHom(P,-)\colon\CondAbl\longrightarrow\CondAbl,\qquad
M\longmapsto\Bigl(S\mapsto\Homop\bigl(P\otimes\mathbb{Z}[S],\,M\bigr)\Bigr),
\end{equation}
whose value on the point is the group of null sequences in $M$. By internal
projectivity, $\iHom(P,-)$ is exact; being a right adjoint (to $P\otimes-$) it
preserves all limits, and it preserves all colimits as well. This exactness is
the engine of everything below.

Consider the endomorphism ``identity minus shift''
\begin{equation}
\label{eq:5:shift}
f=\mathrm{id}-\mathrm{shift}\colon P\longrightarrow P,\qquad
[n]\longmapsto[n]-[n+1].
\end{equation}

\begin{definition}[Solid abelian group]
\label{def:5:solid}
An object $M\in\CondAbl$ is \emph{solid} if the induced map
\begin{equation}
\label{eq:5:solid-def}
f^{\ast}\colon\iHom(P,M)\xrightarrow{\ \sim\ }\iHom(P,M)
\end{equation}
is an isomorphism. Write $\Solid\subset\CondAbl$ for the full subcategory of solid
objects.
\end{definition}

\begin{remark}[The condition is unique summability]
\label{rmk:5:summation-interpretation}
Evaluating at a point, $\iHom(P,M)$ is the space of null sequences
$(m_0,m_1,\dots)$ in $M$, and $f^\ast$ is the difference operator
\begin{equation}
\label{eq:5:difference-operator}
(m_0,m_1,m_2,\dots)\longmapsto(m_0-m_1,\ m_1-m_2,\ \dots).
\end{equation}
Its inverse must be the telescoping/summation operator
\begin{equation}
\label{eq:5:telescope}
(m_0,m_1,m_2,\dots)\longmapsto
\bigl(m_0+m_1+m_2+\cdots,\ m_1+m_2+\cdots,\ \dots\bigr),
\end{equation}
recovering $m_0=(m_0-m_1)+(m_1-m_2)+\cdots$ by telescoping. Thus \eqref{eq:5:solid-def}
says exactly that every null sequence in $M$ has a unique sum, with the tails
themselves forming a null sequence. Imposing the condition on the \emph{internal}
Hom, rather than merely on the naive Hom of null sequences, is a strengthening
that turns out to be much better behaved.
\end{remark}

\subsection{Solid abelian groups form an abelian category}

\begin{proposition}[$\Solid$ is a well-behaved abelian subcategory]
\label{prop:5:solid-abelian}
The subcategory $\Solid\subset\CondAbl$ is an abelian subcategory, stable under
kernels, cokernels, extensions, all limits and all colimits, internal Hom
$\iHom$, and internal Ext $\iExt^{i}$; and it contains $\mathbb{Z}$. In the
terminology to be introduced later, this proposition exhibits an \emph{analytic
ring structure on the integers}.
\end{proposition}

\begin{proof}
Because $P$ is internally projective, $\iHom(P,-)$ is exact and preserves all
limits and colimits; the endomorphism $f$ of \eqref{eq:5:shift} induces a natural
endomorphism $f^\ast$ of $\iHom(P,-)$, and solidity is the assertion that $f^\ast$
is an isomorphism on the object. Every closure property is now inherited from a
corresponding property of the functor $\iHom(P,-)$.

\emph{Kernels, cokernels, extensions, limits, colimits.} Given a map of solid
groups $M\to M'$, apply the exact functor $\iHom(P,-)$: it carries the kernel to
the kernel and the cokernel to the cokernel, and $f^\ast$ acts compatibly, so it
is an isomorphism on both --- the same condition holds for the (co)kernel. For an
extension $0\to M'\to M\to M''\to 0$ with $M',M''$ solid, exactness produces a
short exact sequence of the $\iHom(P,-)$'s on which $f^\ast$ is an isomorphism on
the outer terms, hence on the middle by the five lemma. Stability under all limits
and colimits is likewise immediate, since $\iHom(P,-)$ preserves them and being
solid is a pointwise-isomorphism condition.

\emph{Internal Hom and internal Ext.} Let $M$ be solid and $N\in\CondAbl$
arbitrary; we claim $\iHom(N,M)$ is solid. By tensor--Hom adjunction,
\begin{equation}
\label{eq:5:hom-stability}
\iHom\bigl(P,\iHom(N,M)\bigr)
=\iHom\bigl(P\otimes N,\,M\bigr)
=\iHom\bigl(N,\iHom(P,M)\bigr),
\end{equation}
and $f^\ast$ becomes the map induced by $f$ on the last object, which is an
isomorphism because $M$ is solid; so $f^\ast$ is an isomorphism on
$\iHom(P,\iHom(N,M))$, i.e.\ $\iHom(N,M)$ is solid. Since $\iHom(P,-)$ is exact,
the identical manipulation with $\iExt^i$ in place of $\iHom$ gives
\begin{equation}
\label{eq:5:ext-stability}
\iHom\bigl(P,\iExt^{i}(N,M)\bigr)=\iExt^{i}\bigl(P\otimes N,\,M\bigr)
=\iExt^{i}\bigl(N,\iHom(P,M)\bigr),
\end{equation}
so $\iExt^{i}(N,M)$ is solid as well.

\emph{$\mathbb{Z}$ is solid.} A null sequence in $\mathbb{Z}$ is eventually zero,
so
\begin{equation}
\label{eq:5:hom-P-Z}
\iHom(P,\mathbb{Z})=\bigoplus_{\mathbb{N}}\mathbb{Z},
\end{equation}
and $f=\mathrm{id}-\mathrm{shift}$ is an isomorphism on $\bigoplus_{\mathbb{N}}
\mathbb{Z}$: an eventually-zero sequence is trivially summable. (A little care is
needed to see this internally, again by adjunction.) Hence $\mathbb{Z}$ is solid,
and by the closure properties so is everything built from it.
\end{proof}

\begin{remark}[The point of the argument]
\label{rmk:5:for-free}
Previously, constructing an analytic ring structure was something of an art:
there were many conditions to check and none easy to guarantee, so essentially
only two examples were built by hand --- the solid structure on $\mathbb{Z}$ and
the liquid structure on $\mathbb{R}$ and related rings --- by hard, handcrafted
proofs. Because $P$ is internally projective, imposing the isomorphism condition
\eqref{eq:5:solid-def} for the single endomorphism $f$ makes all of the stability
properties come for free; for any endomorphism of $P$ one could ask such a
condition. What is \emph{not} for free is the existence of a nonzero object
satisfying it, which is exactly the content of \eqref{eq:5:hom-P-Z}.
\end{remark}

The inclusion of a subcategory closed under all limits admits a left adjoint.

\begin{corollary}[Solidification]
\label{cor:5:solidification}
The inclusion $\Solid\subset\CondAbl$ admits a left adjoint, the
\emph{solidification}
\begin{equation}
\label{eq:5:solidification}
(-)^{\solid}\colon\CondAbl\longrightarrow\Solid,\qquad M\longmapsto M^{\solid},
\end{equation}
characterized by
\begin{equation}
\label{eq:5:solidification-adjunction}
\Homop(M,N)=\Homop(M^{\solid},N)\qquad\text{for all }N\in\Solid.
\end{equation}
It is colimit-preserving, and $\Solid$ acquires a unique symmetric monoidal tensor
product $\otimes^{\solid}$ making $(-)^{\solid}$ symmetric monoidal, given
concretely by
\begin{equation}
\label{eq:5:tensor-solid}
M\otimes^{\solid}N:=(M\otimes N)^{\solid}.
\end{equation}
\end{corollary}

\begin{proof}[Sketch]
Existence of $(-)^{\solid}$ is the adjoint functor theorem: the inclusion
preserves all limits, and under mild set-theoretic hypotheses (presentability)
the left adjoint then exists. For the monoidal structure one must check
\begin{equation}
\label{eq:5:tensor-monoidal-check}
(M\otimes N)^{\solid}\xrightarrow{\ \sim\ }\bigl(M^{\solid}\otimes N^{\solid}\bigr)^{\solid},
\end{equation}
which one verifies by mapping into an arbitrary solid $S$: by adjunction and
tensor--Hom,
\[
\Homop(M\otimes N,S)=\Homop\bigl(N,\iHom(M,S)\bigr),
\]
and since $\iHom(M,S)$ is solid (\cref{prop:5:solid-abelian}) one may replace $N$
by $N^{\solid}$; replacing $M$ by $M^{\solid}$ symmetrically gives
$\Homop(M^{\solid}\otimes N^{\solid},S)$. The key input is exactly stability of
$\Solid$ under internal Hom.
\end{proof}

\begin{remark}[A possible right adjoint]
\label{rmk:5:right-adjoint}
Since $\Solid$ is also stable under all colimits, the inclusion preserves
colimits and so, formally, often admits a \emph{right} adjoint as well; Scholze
remarked he had never had occasion to consider it.
\end{remark}

\subsection{The reals are annihilated}

We now record how solidity interacts with the real numbers: not only is
$\mathbb{R}$ not solid, but it admits no nonzero map to any solid group.

\begin{lemma}[$\mathbb{R}^{\solid}=0$]
\label{lem:5:R-solid-zero}
The solidification of the reals vanishes:
\begin{equation}
\label{eq:5:R-solid-zero}
\underline{\mathbb{R}}^{\solid}=0 .
\end{equation}
\end{lemma}

\begin{proof}
Solidification of a ring is a ring (the solid tensor product is symmetric
monoidal), so $\underline{\mathbb{R}}^{\solid}$ is a ring and it suffices to show
$1=0$ in it. Consider the null sequence in $\mathbb{R}$
\begin{equation}
\label{eq:5:R-nullseq}
1,\ \tfrac12,\tfrac12,\ \tfrac14,\tfrac14,\tfrac14,\tfrac14,\ \dots
\qquad\bigl(\tfrac{1}{2^{n}}\text{ repeated } 2^{n}\text{ times}\bigr),
\end{equation}
i.e.\ a map $P\to\underline{\mathbb{R}}$. Composing to
$\underline{\mathbb{R}}\to\underline{\mathbb{R}}^{\solid}$ and using that the target
is solid, the null sequence becomes uniquely summable: there is a unique map from
the summed sequence, intuitively the sequence of tails. Restricting along the
inclusion $\mathbb{Z}\xrightarrow{[0]}P$ of the first term picks out the total sum
\begin{equation}
\label{eq:5:x-total-sum}
x:=1+\tfrac12+\tfrac12+\tfrac14+\tfrac14+\tfrac14+\tfrac14+\cdots\ \in\
\underline{\mathbb{R}}^{\solid}.
\end{equation}
Now
\begin{equation}
\label{eq:5:x-equals-1-plus-x}
x=1+\Bigl(\underbrace{\tfrac12+\tfrac12}+\underbrace{\tfrac14+\tfrac14}
+\underbrace{\tfrac14+\tfrac14}+\cdots\Bigr)
=1+\bigl(1+\tfrac12+\tfrac12+\cdots\bigr)=1+x ,
\end{equation}
where in the second step one is allowed to first sum adjacent pairs --- a
finite-character operation which cannot change the value --- reproducing the same
series. Hence $x=1+x$, so $1=0$ and $\underline{\mathbb{R}}^{\solid}=0$.
\end{proof}

\begin{remark}[Provenance of the argument]
\label{rmk:5:R-argument-history}
Scholze noted that finding an argument for \cref{lem:5:R-solid-zero} took him and
Clausen considerable effort; the divergent-null-sequence argument above he had
worked out only the day before the lecture.
\end{remark}

\begin{corollary}[$\underline{\mathbb{R}}$-modules are killed]
\label{cor:5:R-modules-killed}
If $M\in\CondAbl$ admits an $\underline{\mathbb{R}}$-module structure, then
$M^{\solid}=0$; in fact
\begin{equation}
\label{eq:5:R-module-killed}
\iExt^{i}(M,\,\text{solid})=0\qquad\text{for all }i\ge0 .
\end{equation}
\end{corollary}

\begin{proof}
Resolve the $\underline{\mathbb{R}}$-module $M$ by the bar resolution of free
$\underline{\mathbb{R}}$-modules,
\begin{equation}
\label{eq:5:bar-resolution}
\cdots\longrightarrow M\otimes\underline{\mathbb{R}}\otimes\underline{\mathbb{R}}
\longrightarrow M\otimes\underline{\mathbb{R}}\longrightarrow M\longrightarrow0
\qquad(\text{exact}).
\end{equation}
Solidifying (which is right exact), each term $M\otimes\underline{\mathbb{R}}$
becomes
\[
(M\otimes\underline{\mathbb{R}})^{\solid}
=M^{\solid}\otimes^{\solid}\underline{\mathbb{R}}^{\solid}=0
\]
by \cref{lem:5:R-solid-zero}, so $M^{\solid}$ is a colimit of zeros, i.e.\
$M^{\solid}=0$. Equivalently, the groups $\iExt^{i}(M,\text{solid})$ are modules
over $\underline{\mathbb{R}}^{\solid}=0$, hence vanish. (For the $\iExt$ statement
one uses the derived solidification, or one observes directly that the internal
Ext against a solid object is an $\underline{\mathbb{R}}^{\solid}$-module.)
\end{proof}

\begin{remark}[It is about the topology, not the field]
\label{rmk:5:topology-not-field}
The obstruction is topological, not algebraic. A characteristic-zero field with
its \emph{discrete} topology is perfectly solid; the rationals $\mathbb{Q}$ with
the \emph{subspace} topology from $\mathbb{R}$ are killed, just as $\mathbb{R}$
is. Likewise anything pro-discrete --- an inverse limit of discrete abelian
groups, such as $\mathbb{Z}_p=\varprojlim_n\mathbb{Z}/p^n$ --- is solid.
\end{remark}

\subsection{Computing the solidification of \texorpdfstring{$P$}{P}}

The goal is now to compute $P^{\solid}$. Solidifying $P$ forces in many new
elements: there is a genuine null sequence in $P$, so its sum $1+1+1+\cdots$ (of
the constant sequence) must become defined, and so on; one expects $P^{\solid}$
to be much larger than $P$. The precise statement passes through a bounded
product.

\begin{notation}[Bounded product]
\label{not:5:bounded-product}
Set
\begin{equation}
\label{eq:5:bounded-product-def}
\prod_{\mathbb{N}}^{\mathrm{bdd}}\mathbb{Z}
:=\bigcup_{n\in\mathbb{N}}\ \prod_{\mathbb{N}}\bigl(\mathbb{Z}\cap[-n,n]\bigr)
\ \subseteq\ \prod_{\mathbb{N}}\mathbb{Z},
\end{equation}
the subspace of bounded integer sequences (an increasing union of profinite
subspaces). The standard basis vectors $e_n=(0,\dots,0,1,0,\dots)$ ($1$ in the
$n$-th coordinate) form a null sequence in $\prod_{\mathbb{N}}^{\mathrm{bdd}}
\mathbb{Z}$, because its topology is that of pointwise (coordinatewise)
convergence; this gives a map $P\to\prod_{\mathbb{N}}^{\mathrm{bdd}}\mathbb{Z}$,
$[n]\mapsto e_n$.
\end{notation}

\begin{lemma}[$P^{\solid}$ via the bounded product]
\label{lem:5:P-bounded-product}
The map $[n]\mapsto e_n$ induces an isomorphism
\begin{equation}
\label{eq:5:P-bounded-iso}
P^{\solid}\xrightarrow{\ \sim\ }\Bigl(\prod_{\mathbb{N}}^{\mathrm{bdd}}\mathbb{Z}
\Bigr)^{\solid},
\end{equation}
and more precisely, for every solid $M$ and all $i\ge0$,
\begin{equation}
\label{eq:5:P-bounded-ext}
\iExt^{i}\Bigl(\prod_{\mathbb{N}}^{\mathrm{bdd}}\mathbb{Z},\,M\Bigr)
\xrightarrow{\ \sim\ }\iExt^{i}(P,M).
\end{equation}
\end{lemma}

\begin{proof}[Idea]
Both the identity of $\prod_{\mathbb{N}}^{\mathrm{bdd}}\mathbb{Z}$ and the
composite through $f\otimes\mathrm{id}$ are expressed as sums of null sequences of
endomorphisms. Writing $T=\prod_{\mathbb{N}}^{\mathrm{bdd}}\mathbb{Z}$, a map
$P\otimes T\to X$ is the same as a null sequence of maps $T\to X$ (a map
$P\to\iHom(T,X)$). The evaluation/sum map $P\otimes T\to T$ corresponds to the
null sequence of endomorphisms $p_n\colon T\to T$, ``projection to the $n$-th
coordinate'', whose formal sum is the identity; the composite obtained after
applying $f\otimes\mathrm{id}$ corresponds instead to the tails $q_n$,
``projection to the coordinates $\ge n$'', with $q_n-q_{n+1}=p_n$. Each $p_n$
factors through a single copy of $\mathbb{Z}$, hence through the image of
$P\to T$; this is exactly where the \emph{bounded} product is needed, since for
the full product the coordinates would be unbounded and the factorization would
fail.

Solidifying the resulting square, $f\otimes\mathrm{id}$ becomes an isomorphism
$(f\otimes\mathrm{id})^{\solid}$ (as $f$ solidifies to an isomorphism, by the very
definition of solidity), and the sum map becomes an isomorphism as well
(solidification is monoidal). One reads off that the induced map
$P^{\solid}\to T^{\solid}$ is a split surjection whose splitting composes to the
identity; a short diagram chase upgrades it to an isomorphism (the resulting
idempotent of $P^{\solid}$ is the identity). The identical argument applies to
$\iExt^{i}(-,\text{solid})$.
\end{proof}

Next we replace the bounded product by the full product.

\begin{lemma}[Bounded product versus full product]
\label{lem:5:bounded-to-full}
The inclusion induces an isomorphism
\begin{equation}
\label{eq:5:bounded-full}
\Bigl(\prod_{\mathbb{N}}^{\mathrm{bdd}}\mathbb{Z}\Bigr)^{\solid}
\xrightarrow{\ \sim\ }\Bigl(\prod_{\mathbb{N}}\mathbb{Z}\Bigr)^{\solid}
=\prod_{\mathbb{N}}\mathbb{Z},
\end{equation}
the full product being already solid, and similarly on $\iExt^{i}(-,\text{solid})$.
\end{lemma}

\begin{proof}
Apply $\iExt^{\ast}(-,M)$, $M$ solid, to the short exact sequence
\begin{equation}
\label{eq:5:bounded-full-ses}
0\longrightarrow\prod_{\mathbb{N}}^{\mathrm{bdd}}\mathbb{Z}
\longrightarrow\prod_{\mathbb{N}}\mathbb{Z}
\longrightarrow\prod_{\mathbb{N}}\mathbb{Z}\Big/\prod_{\mathbb{N}}^{\mathrm{bdd}}\mathbb{Z}
\longrightarrow0 .
\end{equation}
It suffices to show the quotient has vanishing $\iExt^{i}(-,\text{solid})$ for all
$i\ge0$. The point is:
\begin{equation}
\label{eq:5:quotient-is-R-module}
\prod_{\mathbb{N}}\mathbb{Z}\Big/\prod_{\mathbb{N}}^{\mathrm{bdd}}\mathbb{Z}
\ \cong\
\prod_{\mathbb{N}}\underline{\mathbb{R}}\Big/\prod_{\mathbb{N}}^{\mathrm{bdd}}
\underline{\mathbb{R}}
\end{equation}
is an $\underline{\mathbb{R}}$-module. Scholze verified this by comparing the
integer and real versions of the short exact sequence
$0\to\prod_{\mathbb{N}}^{\mathrm{bdd}}(-)\to\prod_{\mathbb{N}}(-)\to
\prod_{\mathbb{N}}(-)\big/\prod_{\mathbb{N}}^{\mathrm{bdd}}(-)\to0$ along the
inclusion $\mathbb{Z}\hookrightarrow\underline{\mathbb{R}}$, laid out as a
$3\times3$ diagram with exact rows and columns:
% board 22 (bottom-left, "Claim") @ 01:07 --- boards/board-22.jpg
\[
\begin{tikzcd}[row sep=normal, column sep=small]
0 \arrow[r] & \prod_{\mathbb{N}}^{\mathrm{bdd}}\mathbb{Z} \arrow[r]\arrow[d] & \prod_{\mathbb{N}}^{\mathrm{bdd}}\underline{\mathbb{R}} \arrow[r]\arrow[d] & \prod_{\mathbb{N}}(\mathbb{R}/\mathbb{Z}) \arrow[r]\arrow[d, equal] & 0 \\
0 \arrow[r] & \prod_{\mathbb{N}}\mathbb{Z} \arrow[r]\arrow[d] & \prod_{\mathbb{N}}\underline{\mathbb{R}} \arrow[r]\arrow[d] & \prod_{\mathbb{N}}(\mathbb{R}/\mathbb{Z}) \arrow[r]\arrow[d] & 0 \\
0 \arrow[r] & \prod_{\mathbb{N}}\mathbb{Z}/\prod_{\mathbb{N}}^{\mathrm{bdd}}\mathbb{Z} \arrow[r, "\sim"] & \prod_{\mathbb{N}}\underline{\mathbb{R}}/\prod_{\mathbb{N}}^{\mathrm{bdd}}\underline{\mathbb{R}} \arrow[r] & 0 \arrow[r] & 0
\end{tikzcd}
\]
Because $\mathbb{R}/\mathbb{Z}$ is bounded, the cokernel of
$\mathbb{Z}\hookrightarrow\underline{\mathbb{R}}$ is the \emph{same} object
$\prod_{\mathbb{N}}(\mathbb{R}/\mathbb{Z})$ at both the bounded and the full level
(the ``difference'' Scholze marked on either side), and the comparison map between
these two cokernels is the identity, so the right-hand column is exact with
vanishing bottom term. The nine lemma then forces the bottom row to be exact,
i.e.\ the map \eqref{eq:5:quotient-is-R-module} to be an isomorphism. Being an
$\underline{\mathbb{R}}$-module, it is annihilated by
\cref{cor:5:R-modules-killed}, and the long exact sequence gives \eqref{eq:5:bounded-full}.
\end{proof}

Combining the two lemmas yields the fundamental computation.

\begin{corollary}[$P^{\solid}$ is a product]
\label{cor:5:P-solid}
There is a canonical isomorphism
\begin{equation}
\label{eq:5:P-solid}
P^{\solid}\xrightarrow{\ \sim\ }\prod_{\mathbb{N}}\mathbb{Z},
\end{equation}
and for all solid $M$, $\iExt^{i}(P,M)\xleftarrow{\ \sim\ }\iExt^{i}_{\mathrm{CondAb}}
\bigl(\prod_{\mathbb{N}}\mathbb{Z},M\bigr)$. In particular, for discrete $M$,
\begin{equation}
\label{eq:5:prod-ext}
\Ext^{i}\Bigl(\prod_{\mathbb{N}}\mathbb{Z},\,M\Bigr)=
\begin{cases}
\bigoplus_{\mathbb{N}}M & i=0,\\
0 & i>0,
\end{cases}
\end{equation}
so that $\prod_{\mathbb{N}}\mathbb{Z}$ is a compact projective object of $\Solid$.
\end{corollary}

\begin{remark}[An independent route to \cref{cor:4:solid-generator}]
\label{rmk:5:independent-of-lec4}
Equation \eqref{eq:5:prod-ext} is precisely \cref{cor:4:solid-generator}, the
vanishing of the higher $\Ext$-groups out of $\prod_{\mathbb{N}}\mathbb{Z}$. In
\cref{lec:4} it was a hard input, proved by embedding into products of
$\underline{\mathbb{R}}$; here it comes out as an \emph{output} of the light
presentation of solidity, an independent argument. (The proof of
\cref{lem:5:bounded-to-full} does use, through
\cref{cor:5:R-modules-killed,lem:5:R-solid-zero}, that products of
$\underline{\mathbb{R}}$ are annihilated.)
\end{remark}

The completed tensor product now does what was wanted in
\cref{rmk:5:tensor-motivation}.

\begin{corollary}[The completed tensor square]
\label{cor:5:tensor-square}
There is a canonical isomorphism
\begin{equation}
\label{eq:5:tensor-product-Z}
\prod_{\mathbb{N}}\mathbb{Z}\ \otimes^{\solid}\ \prod_{\mathbb{N}}\mathbb{Z}
\ \cong\ \prod_{\mathbb{N}\times\mathbb{N}}\mathbb{Z},
\end{equation}
coming from $\prod_{\mathbb{N}}\mathbb{Z}\otimes^{\solid}\prod_{\mathbb{N}}\mathbb{Z}
=(P\otimes P)^{\solid}$ (monoidality of solidification) and the two-variable
analogue of \cref{cor:5:P-solid}. Identifying
$\prod_{\mathbb{N}}\mathbb{Z}$ with the power-series ring in one variable (measures,
cf.\ \cref{rmk:1:MRN-ring}), this reads
\begin{equation}
\label{eq:5:power-series-tensor}
\mathbb{Z}[\![T]\!]\ \otimes^{\solid}\ \mathbb{Z}[\![U]\!]
\ \cong\ \mathbb{Z}[\![T,U]\!],
\end{equation}
so the completed tensor product of two power-series rings is the power-series ring
in both variables --- exactly the reasonable answer sought at the outset.
\end{corollary}

\subsection{The solidification of a general free group}

The computation extends from the convergent sequence to an arbitrary profinite
set, and the resulting formula is the one that \emph{defined} solidification in
the original approach.

\begin{proposition}[{Solidification of $\mathbb{Z}[S]$}]
\label{prop:5:general-S}
Let $S$ be any infinite light profinite set, written $S=\varprojlim_{n}S_{n}$
with the $S_{n}$ finite and surjective transition maps. Then there is a map
$g\colon P\to\mathbb{Z}[S]$ inducing an isomorphism
\begin{equation}
\label{eq:5:ZS-solid}
\mathbb{Z}[S]^{\solid}\ \cong\ \prod_{\mathbb{N}}\mathbb{Z},
\end{equation}
and canonically
\begin{equation}
\label{eq:5:ZS-limit}
\mathbb{Z}[S]^{\solid}\ \xrightarrow{\ \sim\ }\ \varprojlim_{n}\mathbb{Z}[S_{n}],
\end{equation}
compatibly with $\iExt^{i}(-,\text{solid})$.
\end{proposition}

\begin{proof}[Sketch]
Let $\pi_{n}\colon S\to S_{n}$ be the projection. Choose sections inductively: a
section $i_{0}\colon S_{0}\hookrightarrow S$ of $\pi_0$; then, on the elements of
$S_{1}$ not in the image of $i_{0}$, a further section $S_{1}\setminus S_{0}\to S$;
gluing gives $i_{1}\colon S_{1}\to S$ (equal to $i_0\pi_0$-lift on old elements,
new on the rest); and so on.
% board 28 @ 01:34:23 — /home/deven/documents/coding/vms/gpu/home/notetaker/output/peter-scholze-524-analytic-stacks/boards/board-28.jpg
\[
\begin{tikzcd}[row sep=large, column sep=large]
& S \arrow[ld, "\pi_n"'] & \\
S_n \arrow[r, two heads] \arrow[ru, bend right=20, dashed, "i_n"'] & S_{n-1} \arrow[r, two heads] & \cdots \arrow[r, two heads] & S_0
\end{tikzcd}
\]
The union of the images is a countable subset of $S$, and one enumerates
\begin{equation}
\label{eq:5:enumerate}
\mathbb{N}=S_{0}\sqcup(S_{1}\setminus S_{0})\sqcup(S_{2}\setminus S_{1})\sqcup\cdots .
\end{equation}
Define $g\colon P\to\mathbb{Z}[S]$ on this enumeration: on $S_{0}$ by $i_{0}$, and
on $S_{n}\setminus S_{n-1}$ by the difference of the maps induced by $i_{n}$ and by
$S_{n}\setminus S_{n-1}\to S_{n-1}\xrightarrow{i_{n-1}}S$. As in
\cref{lem:5:P-bounded-product}, the identity endomorphism of $\mathbb{Z}[S]$ is the
sum of the null sequence of endomorphisms
\begin{equation}
\label{eq:5:endo-nullseq}
\mathrm{id},\ \mathrm{id}-i_{0}\pi_{0},\ \mathrm{id}-i_{1}\pi_{1},\ \dots,\
\mathrm{id}-i_{n}\pi_{n},\ \dots\,,
\end{equation}
equivalently of the differences $i_{0}\pi_{0},\ i_{1}\pi_{1}-i_{0}\pi_{0},\dots$,
each of which factors through the image of $g$ (this uses that $S$ is a
\emph{light} profinite set, so that it is exhausted through countably many finite
quotients). Solidifying as before produces the isomorphism.
\end{proof}

\begin{remark}[The formula that used to be the definition]
\label{rmk:5:was-definition}
In the original set-up of the solid theory
\cite{clausen2019condensed}, the formula $\mathbb{Z}[S]^{\solid}=\varprojlim_{n}
\mathbb{Z}[S_{n}]$ of \eqref{eq:5:ZS-limit} was the \emph{starting point}: one
declared by hand what the solidification of each generator $\mathbb{Z}[S]$ should
be, taking the naive inverse limit, and then checked by hand that this yields a
valid theory --- a much harder route. In the present light presentation, solidity
is defined first (\cref{def:5:solid}) and \eqref{eq:5:ZS-limit} becomes a
computation.
\end{remark}

\subsection{Structure theorem and integration}

\begin{theorem}[Structure of $\Solid$]
\label{thm:5:structure}
The subcategory $\Solid\subset\CondAbl$ is abelian and stable under all limits and
colimits, and:
\begin{enumerate}
\item it has a single compact projective generator $\prod_{\mathbb{N}}\mathbb{Z}$,
which is internally projective and satisfies
\begin{equation}
\label{eq:5:generator-tensor}
\prod_{\mathbb{N}}\mathbb{Z}\ \otimes^{\solid}\ \prod_{\mathbb{N}}\mathbb{Z}
\ \cong\ \prod_{\mathbb{N}\times\mathbb{N}}\mathbb{Z};
\end{equation}
\item an object $M\in\CondAbl$ is solid if and only if, for every $S\in
\ProFin_{\mathbb{N}}(\mathrm{Fin})$ (with $S=\varprojlim_n S_n$) and every map
$g\colon S\to M$, there is a unique extension of $g$ along
$S\to\mathbb{Z}[S]^{\solid}=\varprojlim_{n}\mathbb{Z}[S_{n}]$,
\begin{equation}
\label{eq:5:point-characterization}
\mathbb{Z}[S]^{\solid}=\varprojlim_{n}\mathbb{Z}[S_{n}]\dashrightarrow M .
\end{equation}
\end{enumerate}
\end{theorem}

The characterization \eqref{eq:5:point-characterization} was the \emph{original}
definition of solidity in \cite{clausen2019condensed}; note it involves only
ordinary maps out of $S$, not internal Hom, in contrast to \cref{def:5:solid}.

\begin{proof}
Only the ``$\Leftarrow$'' direction of (2) is new. Present $M$ in $\CondAbl$ as a
cokernel of free groups,
\begin{equation}
\label{eq:5:presentation}
\bigoplus_{j}\mathbb{Z}[S_{j}]\longrightarrow\bigoplus_{i}\mathbb{Z}[S_{i}]
\longrightarrow M\longrightarrow0 .
\end{equation}
The hypothesis of (2) says that every map $\mathbb{Z}[S_{i}]\to M$ extends
uniquely to $\mathbb{Z}[S_{i}]^{\solid}$, and likewise for the $S_{j}$, so the
presentation solidifies to
\[
\bigoplus_{j}\mathbb{Z}[S_{j}]^{\solid}\longrightarrow
\bigoplus_{i}\mathbb{Z}[S_{i}]^{\solid}\dashrightarrow M ,
\]
exhibiting $M$ as a cokernel of solid objects; since $\Solid$ is closed under
cokernels (\cref{prop:5:solid-abelian}), $M$ is solid. Equivalently, the universal
property makes $M$ a \emph{retract} of its solidification $M^{\solid}$ (the
composite $M\to M^{\solid}\to M$ being the identity, and a comparison map going
back), and retracts of solid objects are solid. The forward direction is clear,
since a map $S\to M$ is a map $\mathbb{Z}[S]\to M$, which factors uniquely through
the solidification when $M$ is solid.
\end{proof}

\begin{remark}[Solidity as integration against measures]
\label{rmk:5:measures}
The philosophical content of the condition is integration. One has
\begin{equation}
\label{eq:5:measures}
\mathbb{Z}[S]^{\solid}=\iHom\bigl(\Cont(S,\mathbb{Z}),\,\mathbb{Z}\bigr),
\end{equation}
because $\Cont(S,\mathbb{Z})$ is the filtered colimit of the functions on the
finite quotients $S_n$, and dualizing gives the inverse limit
\eqref{eq:5:ZS-limit}; so $\mathbb{Z}[S]^{\solid}$ is the group of
\emph{$\mathbb{Z}$-valued measures} on $S$. Thus, by \eqref{eq:5:point-characterization},
$M$ is solid precisely when, given a map $g\colon S\to M$ from a profinite set and
a $\mathbb{Z}$-valued measure $\mu$ on $S$, one can form the integral
\begin{equation}
\label{eq:5:integration}
\int_{S}g\,\mathrm{d}\mu\ \in\ M .
\end{equation}
Note this final characterization speaks only of ordinary Hom, whereas the
null-sequence definition \cref{def:5:solid} was phrased on the internal Hom.
\end{remark}

\begin{remark}[Lightness is essential]
\label{rmk:5:lightness-essential}
The characterization by null sequences is genuinely a light phenomenon. Were one
to attempt the same definition for \emph{all} condensed abelian groups, the
null-sequence condition would be too weak: it only constrains countable data,
whereas a general profinite set is not exhausted by countably many finite
quotients. It is exactly the fact that a light profinite set \emph{is} exhausted
through countably many finite quotients that makes \cref{prop:5:general-S} work.
\end{remark}

\begin{remark}[A synthetic description; and the liquid question]
\label{rmk:5:synthetic-and-liquid}
Having a single compact projective generator $\prod_{\mathbb{N}}\mathbb{Z}$ means
$\Solid$ is equivalent to modules over its endomorphism ring, so one can describe
$\Solid$ purely algebraically, with no mention of profinite sets. The morphisms of
the generator are
\begin{equation}
\label{eq:5:endomorphisms}
\Homop\Bigl(\prod_{\mathbb{N}}\mathbb{Z},\prod_{\mathbb{N}}\mathbb{Z}\Bigr)
=\prod_{\mathbb{N}}\bigoplus_{\mathbb{N}}\mathbb{Z},
\end{equation}
i.e.\ infinite matrices each of whose rows is eventually zero; Scholze deferred
this synthetic description to the next lecture. Asked whether the liquid
(archimedean) theory of vector spaces admits an analogous null-sequence-style
characterization, he said he expected it should be possible and that he and
Clausen were currently working out the details.
\end{remark}
